% =============================================================================
%  المذكرة البيداغوجية رقم 01
%  المجال 1: بيئة التعامل مع الحاسوب — الوحدة 1: تقنية المعلومات (IT)
%  المورد المستهدف: مفهوم تكنولوجيا الإعلام والاتصال (TIC)، مفهوم المعلوماتية،
%                   والفرق بين البيانات والمعلومات
%  السنة الدراسية: 2026 / 2027
% =============================================================================
%  البناء (من هذا المجلد):
%    xelatex -interaction=nonstopmode -halt-on-error -output-directory=build lesson-note.tex
% =============================================================================

\documentclass[11pt, a4paper]{article}

% ── المشترك (التمهيد + التصميم + الماكروهات) ────────────────────────────────
\input{"../../shared/common.tex"}

% ══════════════════════════════════════════════════════════════════════════════
\begin{document}

% ── الصفحة 1: الترويسة + البطاقة الفنية ─────────────────────────────────────
\officialheader

\cardtitle{مذكرة بيداغوجية رقم: 01}

{\gridsetup
\centering
\begin{tabularx}{0.85\textwidth}{@{}|c|X|@{}}
  \hline
  \cardrow{المؤسسة}{ثانوية الشهيد حكومي العيد أدرار — أدرار}
  \cardrow{المادة}{المعلوماتية}
  \cardrow{المستوى والشعبة}{السنة الأولى ثانوي (جذع مشترك علوم وتكنولوجيا / جذع مشترك آداب وفلسفة)}
  \cardrow{التاريخ}{28 سبتمبر 2026}
  \cardrow{المدة}{02 ساعتان}
  \cardrow{المجال التعلمي}{بيئة التعامل مع الحاسوب}
  \cardrow{الوحدة التعلمية}{تقنية المعلومات (IT)}
  \cardrow{المورد المستهدف}{مفهوم تكنولوجيا الإعلام والاتصال (TIC)، مفهوم المعلوماتية، والفرق بين البيانات والمعلومات}
  \cardrow{الكفاءة المستهدفة}{اكتساب المفاهيم القاعدية الأساسية لتقنية المعلومات والتمكّن من التمييز بين المادة الخام (البيانات) والنتائج المعالجة (المعلومات).}
  \cardrow{الأهداف التعليمية}{يتعرّف التلميذ على مفاهيم التقنية، التكنولوجيا، الإعلام، والاتصال.\newline
    يتعرّف على مفهوم تكنولوجيا الإعلام والاتصال (TIC) ومكوناتها المادية والشبكية.\newline
    يميّز بين البيانات والمعلومات وعملية المعالجة الآلية.\newline
    يتعرّف على مفهوم المعلوماتية (Informatique) وأصل الكلمة.
    يتعرّف على مراحل تطور تقنية المعلومات (تطور الأجهزة والبرامج).}
  \cardrow{المكتسبات القبلية}{المفاهيم العامة الأولية حول الحاسوب والمعلوماتية.}
  \cardrow{الوسائل والدعامات}{العارض الرقمي (Data Show)، الكتاب المدرسي، ومخططات توضيحية لدورة معالجة البيانات.}
\end{tabularx}\par
}\vspace{8pt}

% ── الصفحة 2: مراحل سير الوضعية التعلمية ────────────────────────────────────
\lessonbreak{مراحل سير الوضعية التعلمية}

\stageheader
\stagecontent
  {\badge{1}\\[2pt] مرحلة الانطلاق\\\footnotesize (15 دقيقة)}
  {التذكير بالمكتسبات القبلية وطرح الوضعية المشكلة.\newline
   \textbf{الوضعية المشكلة:} في حياتنا اليومية نستعمل مصطلحات مثل «المعلوماتية»، «الإعلام الآلي»، «التكنولوجيا»، و«الإعلام والاتصال»، ولكلٍّ منها مجالات استعمال معروفة: إنجاز الدروس والواجبات ومتابعة السلع والفواتير في المحلات (المعلوماتية والإعلام الآلي)، والبيوت الذكية والروبوتات والزراعة الآلية والطب عن بُعد (التكنولوجيا)، والإذاعة والبثّ المباشر على منصات التواصل (الإعلام والاتصال). والقاسم المشترك بينها جميعًا أنها تُيسِّر حياة الإنسان.\newline
   \textbf{الإشكالية:} هل تعبّر هذه الكلمات عن نفس المعنى؟ نلاحظ أنّ «المعلوماتية» و«الإعلام الآلي» يدلّان على المعنى نفسه، أمّا «التكنولوجيا» و«الإعلام والاتصال» فمفاهيم مختلفة. فما مفهوم تقنية المعلومات؟ وما مراحل تطورها؟}
  {يذكّر التلاميذ بالمكتسبات القبلية، ويعرض المصطلحات الأربعة مع مجالات استعمال كلٍّ منها، ويوجّه المناقشة نحو الترادف الجزئي بين المصطلحات.}
  {يقدّم التلاميذ أمثلة من تجربتهم اليومية لكل مجال استعمال، ويستنتجون أنّ بعض المصطلحات مترادف وبعضها مختلف.}
  {وضعية مشكلة ومناقشة جماعية}
  {تقويم تشخيصي}

\stagecontent
  {\badge{2}\\[2pt] بناء التعلمات\\\footnotesize (65 دقيقة)}
  {\textbf{المحاور والمعارف المستهدفة:}\newline
   التفكيك المفهومي للتقنية والتكنولوجيا، مفهوم TIC، البيانات والمعلومات والمعالجة، مفهوم المعلوماتية، ثمّ مراحل تطور تقنية المعلومات (تطور أجهزة المعالجة، وتطور البرامج: أنظمة التشغيل والتطبيقات).}
  {يطرح الأستاذ أسئلة لاستكشاف التمييز بين المهارة الفردية والعلم التطبيقي، ويوجّه التلاميذ للربط بين مفهومي الإعلام والاتصال، ويعرض مجموعة بيانات غير مرتبة على السبورة ويطلب صياغتها في جملة مفيدة، ثم يشرح التركيب اللغوي للمصطلح العلمي Informatique، ويعرض على العارض الرقمي مراحل تطور تقنية المعلومات انطلاقًا من حاسوب بحجم غرفة إلى الهاتف الذكي، مع ذكر أشهر أنظمة التشغيل وتطبيقاتها.}
  {يستنتج التلاميذ أصل الكلمة وتطبيقاتها، ويحلّلون مكونات تكنولوجيا الإعلام والاتصال، ثم يرتّبون مراحل تطور الأجهزة والبرامج ويذكرون أمثلة من حولهم.}
  {أسئلة استكشافية وعرض شارح}
  {تقويم تكويني}

\stagecontent
  {\badge{3}\\[2pt] الحصيلة والتقويم\\\footnotesize (40 دقيقة)}
  {صياغة الخلاصة وتثبيت المفاهيم، ثم إنجاز النشاط التطبيقي (تقويم مرحلي).}
  {يقدّم الأستاذ ثلاثة تمارين تطبيقيين للتأكد من فهم دورة معالجة البيانات ومراحل تطور تقنية المعلومات.}
  {ينجز التلاميذ التمرينات الثلاثة ويقدّمون إجاباتهم ومناقشة النتائج.}
  {تمارين تطبيقي وتقويم تتابعي}
  {تقويم ختامي}

% ── الصفحة 3: الصياغة والتأسيس المعرفي ──────────────────────────────────────
\lessonbreak{الصياغة والتأسيس المعرفي}

\begin{rememberbox}[المصطلحات ومجالات استعمالها]
  \begin{itemize}[label=\textbullet, leftmargin=1.1em, itemsep=2pt, topsep=1pt, parsep=0pt]
    \item \textbf{المعلوماتية / الإعلام الآلي:} إنجاز الدروس والواجبات، متابعة السلع والفواتير في المحلات، النظام المعلوماتي لوزارة التربية (معالجة الطلبات، المسابقات، تسجيل التلاميذ لشهادة التعليم المتوسط).
    \item \textbf{التكنولوجيا:} البيوت الذكية (مثل Siri)، الروبوتات، الزراعة الآلية، الطبّ عن بُعد.
    \item \textbf{الإعلام والاتصال:} الإذاعة والبثّ (عتاد وميكروفون وكاميرا وتطبيقات المونتاج)، والبثوث المباشرة على منصات التواصل.
    \item \textbf{القاسم المشترك:} جميعها تُيسِّر حياة الإنسان.
  \end{itemize}
\end{rememberbox}

\begin{rememberbox}[1. التفكيك المفهومي للتقنية والتكنولوجيا]
  \begin{itemize}[label=\textbullet, leftmargin=1.1em, itemsep=2pt, topsep=1pt, parsep=0pt]
    \item \textbf{التقنية (Technique):} المهارة والقدرة على ممارسة الأعمال والأنشطة لإنجاز عمل ما.
    \item \textbf{التكنولوجيا (Technologie):} كلمة يونانية مركّبة من مقطعين (Techno: مهارة فنية + Logos: علم)، وتعني علم المهارات التطبيقية والاستغلال المنطقي للمعارف العلمية.
  \end{itemize}
\end{rememberbox}

\begin{rememberbox}[2. مفهوم تكنولوجيا الإعلام والاتصال (TIC)]
  \begin{itemize}[label=\textbullet, leftmargin=1.1em, itemsep=2pt, topsep=1pt, parsep=0pt]
    \item \textbf{التعريف:} التلاقي والتزاوج بين عتاد الحواسيب والبرمجيات وشبكات الاتصال لجمع البيانات وتخزينها ومعالجتها ونقلها.
    \item \textbf{المكونات الأساسية:} تقنيات المعالجة والتخزين (أجهزة وبرامج)، وشبكات الاتصال السلكية واللاسلكية.
  \end{itemize}
\end{rememberbox}

\begin{rememberbox}[3. البيانات والمعلومات والمعالجة]
  \begin{itemize}[label=\textbullet, leftmargin=1.1em, itemsep=2pt, topsep=1pt, parsep=0pt]
    \item \textbf{البيانات (Data / Données):} القيم والحقائق الأولية الخام (أرقام، حروف، صور) قبل تنظيمها.
    \item \textbf{المعالجة (Processing / Traitement):} العمليات الحسابية أو المنطقية المطبّقة على البيانات.
    \item \textbf{المعلومات (Informations):} البيانات التي تمّت معالجتها وأصبحت ذات معنى وسياق محدّد.
  \end{itemize}
\end{rememberbox}

\begin{rememberbox}[4. مفهوم المعلوماتية (Informatique)]
  \begin{itemize}[label=\textbullet, leftmargin=1.1em, itemsep=2pt, topsep=1pt, parsep=0pt]
    \item \textbf{أصل الكلمة:} كلمة مركّبة من كلمتين فرنسيتين هما (Information: معلومات) + (Automatique: آلية).
    \item \textbf{التعريف:} العلم الذي يعالج المعلومات بطريقة آلية باستخدام جهاز الحاسوب وفق برامج محدّدة.
  \end{itemize}
\end{rememberbox}

% ── مراحل تطور تقنية المعلومات ──────────────────────────────────────────────
\lessonbreak{مراحل تطور تقنية المعلومات}

\begin{rememberbox}[أ. تطور أجهزة المعالجة]
  \begin{itemize}[label=\textbullet, leftmargin=1.1em, itemsep=2pt, topsep=1pt, parsep=0pt]
    \item تطوّرت أجهزة المعالجة من حاسوب يشغل حجم غرفة كاملة، إلى حاسوب محمول، ثمّ جهاز لوحي، ثمّ هاتف ذكي.
    \item ويشمل التطوّر أيضًا مكوّنات الحاسوب، مثل الرامات وبطاقة الشاشة.
  \end{itemize}
\end{rememberbox}

\begin{rememberbox}[ب. تطور البرامج: أنظمة التشغيل]
  \begin{itemize}[label=\textbullet, leftmargin=1.1em, itemsep=2pt, topsep=1pt, parsep=0pt]
    \item \textbf{التعريف:} تسمح باستغلال الحاسوب وتشغيل بقيّة البرامج الأخرى.
    \item \textbf{أنظمة التشغيل الثلاثة:} Windows وLinux وmacOS.
    \item \textbf{مسار تطوّر Windows:} MS-DOS ثمّ 1.0 و2.0 و3.0، ثمّ 95 و98، فـ XP وVista، ثمّ 7 و8، ثمّ 10 و11.
    \item \textbf{Linux (أشهر التوزيعات):} Ubuntu وArch Linux وRed Hat وFedora.
  \end{itemize}
\end{rememberbox}

\begin{rememberbox}[ج. تطور التطبيقات]
  \begin{itemize}[label=\textbullet, leftmargin=1.1em, itemsep=2pt, topsep=1pt, parsep=0pt]
    \item تشمل البرامج التطبيقية: تطبيقات Adobe، والمتصفحات، والألعاب، وبرامج Word وExcel وPowerPoint.
    \item يُعدّ كلُّ تحديثٍ لهذه التطبيقات مرحلةً من مراحل تطور تقنية المعلومات.
  \end{itemize}
\end{rememberbox}

% ── الصفحة 4: النشاط التطبيقي ───────────────────────────────────────────────
\lessonbreak{النشاط التطبيقي (تقويم مرحلي)}

\textbf{\color{primary}التمرين الأول:} صنّف العناصر التالية في جدول ينقسم إلى (بيانات / معلومات):
\begin{enumerate}[label=\arabic*., leftmargin=1.4em, itemsep=2pt, topsep=2pt, parsep=0pt]
  \item أرقام النقاط المفردة للتلاميذ دون أسماء (12، 08، 16).
  \item كشف النقاط الفصلي وملاحظة «يرتقي».
  \item اسم التلميذ، تاريخ الميلاد، وعنوان السكن في استمارة التسجيل.
  \item القائمة الاسمية للتلاميذ الناجحين مرتّبة حسب المعدل.
\end{enumerate}

{\gridsetup
\centering
\begin{tabularx}{0.85\textwidth}{@{}|>{\centering\arraybackslash}X|>{\centering\arraybackslash}X|@{}}
  \hline
  \rowcolor{primary}
  \thead{بيانات (Data)} & \thead{معلومات (Informations)} \\
  \hline
  \answerline & \answerline \\
  \answerline & \answerline \\
  \answerline & \answerline \\
  \hline
\end{tabularx}\par
}\vspace{8pt}

\textbf{\color{primary}التمرين الثاني:} بيّن في مخطط سهمي بسيط كيف تتكامل البيانات والمعالجة للحصول على المعلومات.

\vspace{4pt}
\begin{center}
\begin{tikzpicture}[
  node distance=2.6cm,
  box/.style={draw=primary, fill=lightbg, thick, rounded corners=3pt,
              minimum width=3.6cm, minimum height=1.1cm, align=center},
  ar/.style={-{Latex[length=3mm]}, thick, primary}
]
  \node[box] (data) {بيانات أولية\\\footnotesize (Data)};
  \node[box, right=of data] (proc) {معالجة آلية\\\footnotesize (Processing)};
  \node[box, right=of proc] (info) {معلومات مفيدة\\\footnotesize (Information)};
  \draw[ar] (data) -- (proc);
  \draw[ar] (proc) -- (info);
\end{tikzpicture}
\end{center}

\vspace{6pt}
\textbf{\color{primary}التمرين الثالث:} صنّف العناصر التالية في جدول ثلاثي ينقسم إلى (أجهزة المعالجة / أنظمة التشغيل / التطبيقات):
\begin{enumerate}[label=\arabic*., leftmargin=1.4em, itemsep=2pt, topsep=2pt, parsep=0pt]
  \item بطاقة الشاشة.
  \item نظام Ubuntu.
  \item برنامج Word.
  \item هاتف ذكي.
  \item نظام MS-DOS.
  \item تطبيقات Adobe.
  \item الرامات.
\end{enumerate}

{\gridsetup
\centering
\begin{tabularx}{0.85\textwidth}{@{}|>{\centering\arraybackslash}X|>{\centering\arraybackslash}X|>{\centering\arraybackslash}X|@{}}
  \hline
  \rowcolor{primary}
  \thead{أجهزة المعالجة} & \thead{أنظمة التشغيل} & \thead{التطبيقات} \\
  \hline
  \answerline & \answerline & \answerline \\
  \answerline & \answerline & \answerline \\
  \hline
\end{tabularx}\par
}\vspace{8pt}

\begin{rememberbox}[الخلاصة المدوّنة]
  \begin{center}
    بيانات أوليّة (Data) $\xrightarrow{\ \ \text{معالجة آلية (Processing)}\ \ }$ معلومات مفيدة (Information)
  \end{center}
\end{rememberbox}

\end{document}
